\section{The Hard Cases}
\label{sec:hard}

Each case below is stated, run through the rules, and marked \emph{solved}, \emph{solved with a
cost}, or \emph{open}.

\subsection{Closures that read a captured list, and closures called many times}
\label{sec:hard-closures}

\begin{Verbatim}
total = List.foldl xs 0 (λx acc → acc + List.length ys)    -- reads ys, called n times
zs = List.set ys 0 7                                       -- then edits ys
\end{Verbatim}
The lambda's summary is $\mathit{cap}\ \mathit{ys} = \borrowed$ and $\mathit{once} = \mathit{no}$.
\code{foldl}'s class for its callback is $\mathit{calls} = \mathit{many}$, $\mathit{escapes} =
\mathit{no}$, and a closure that is not $\mathit{once}$ may be called many times. The closure holds
\code{ys} while it is live, which is during the \code{foldl} call only, so at \code{List.set ys 0 7} the
holder is dead. \emph{Solved.} The reason is that beni's lambdas are in front of the checker. In
type-based systems the capture counts as a use of the capture's \emph{type}: in Wansbrough and Peyton
Jones's usage analysis, ``all free variables of an abstraction have at least the usage of the
abstraction itself'' \citep[\S6.1]{wansbrough1999once}, and effect- and capture-based systems close
the case only with declarations \citep{deng2025capture,xu2026capybara,radanne2020kindly}.

\begin{Verbatim}
ys2 = List.map xs (λx → List.push acc x)                    -- consumes a capture, called n times
\end{Verbatim}
Here $\mathit{cap}\ \mathit{acc} = \consumed$ and $\mathit{once} = \mathit{yes}$, while \code{map}'s
class says $\mathit{calls} = \mathit{many}$. The edge $\mathit{once} \Rightarrow \mathit{calls} \le 1$
fails, and the program is \emph{rejected}, correctly: the first call writes \code{acc} and returns it,
the second writes the same array, and the first element of the result would change under the second
call. The message names the lambda, the capture and \code{map}'s call count, and offers a fold
(\S\ref{sec:errors-examples}). This is OxCaml's lock rule for closures, under which a closure that
consumes a captured value must be \code{once} \citep[\S3.3]{lorenzen2024oxidizing}, and Futhark's
refusal to let the function of a \code{map} consume an array bound outside it, since ``that would
imply the array is consumed once for every iteration'' \citep[\S3.3]{henriksen2017futhark}. A closure
that consumes a capture and \emph{is} called at most once---\code{Task.scope (λscope → … List.push acc
x …)}, \code{Result.andThen r (λx → List.push acc x)}---passes, because those callees' classes say
$\mathit{calls} \le 1$. \emph{Solved}, with a cost: a $\mathit{once}$ closure passed to a
$\mathit{many}$ position is an error even when the program would in fact call it once, for example a
\code{List.map} over a one-element list.

\subsection{Pipelines, identity and composition}
\label{sec:hard-pipelines}

The pipe is syntax: \code{xs ▷ List.filter p ▷ List.set 0 v} \emph{is}
\code{List.set (List.filter xs p) 0 v}, and a placeholder is a written lambda, so there is no plumbing
function in which to lose precision. \code{identity x = x} is inferred with $\mathit{res} = \{\pi_1\}$;
a user-written \code{apply f x = f x} with $\mathit{res}$ equal to its callback's result class; a
composition \code{λx → g (f x)} likewise. Futhark's alias analysis is lossy exactly here and its
authors proposed parametricity as the fix \citep{henriksen2026aliasing}; we analyse the body and let
the class carry the result's provenance. \emph{Solved.}

The residual cost lies in the identity promises of beni's current list library: \code{filter} that
keeps every element returns its input itself, as does \code{map} when every result is \code{===} its
element, and so on. Under the rule such a function has $\mathit{res} = \{\fresh, \pi_1\}$---the result
\emph{may} be the input---so \code{List.set (List.filter xs p) 0 v} demands \code{xs} unique too. In a
pipeline \code{xs} is usually dead afterwards and this costs nothing; where it is not, the rejection
names the promise. We propose dropping the promises that alias (\S\ref{sec:disc-options}).

\subsection{Generic code}
\label{sec:hard-generic}

The well-known \code{eq} and \code{compare} are read-only by contract; a user method's summary is a
class bound at each call site (\S\ref{sec:inf-generic}). A generic body cannot write a value of type
\code{a} in place, which is the right answer. \emph{Solved.}

\subsection{Lists in containers}
\label{sec:hard-containers}

\paragraph{A record.} In \code{\{ m | rows = List.push m.rows x \}} the operand is $(m, .\mathit{rows})$.
A record update shares every \emph{other} field's value between the old and the new record, but not
\code{rows}, which the push's result replaces. After the update \code{m} is dead unless read later,
and a later \code{m.name} reads a path that does not meet $.\mathit{rows}$. \emph{Solved}: a model
threaded through helpers that each touch one field is accepted, because $\mathit{use}$ and $\live$ are
per path. Tuples and constructors are the same, through $.i$ and $C\#i$.

\paragraph{A list of lists.} In \code{List.update rows i (λr → \{ r | tags = List.push r.tags t \})} the
lambda's parameter is an element of \code{rows}. Two slots of \code{rows} may hold the same record, in
which case a push through slot $i$ would be seen at the other slot. The derivation therefore needs two
things: \code{List.update} must be a consuming writer that overwrites slot $i$ with the callback's
result---a destructive read, so the slot's old holder is dead during the callback---and \code{rows} must
be unique and exclusive (\S\ref{sec:model-excl}). Then \code{r} is unique inside the callback,
\code{r.tags}'s only holder is \code{r}, and the push is in place. If \code{rows} is unique but not
exclusive---its elements came from a \code{foreign} that declares nothing, or from
\code{List.repeat}---the program is rejected as a limit, with the fix \code{List.copy r.tags}. With
\code{map} in place of \code{update}, \code{map} is not a writer and supplies a borrowed element: the old
\code{rows} holds every element until \code{map} has built its result, so the push is rejected as real
sharing (unless \code{map} joins the demand set, \S\ref{sec:disc-options}). \emph{Solved with a cost}:
a nested edit needs exclusivity, which the program's own construction of the container must establish.

\paragraph{A dictionary.} beni's \code{Dict} is a balanced tree written in beni. A lookup
\code{Dict.get d k} is a non-destructive read: its result aliases a value inside \code{d}, and \code{d}'s
values stop being exclusive while that result is live. \code{Dict.insert d k v} stores \code{v} (shared),
rebuilds a path and shares the rest with \code{d}. \code{Dict.update d k f} takes the node's value,
calls \code{f} and rebuilds the node with the result---a destructive read, whose callback's argument is
unique when \code{d} is consumed at the call and its values are exclusive. Huisinga's motivating
example for Lean \citep[\S2.1]{huisinga2023lean}, in which a group is looked up, pushed to and
re-inserted, is \emph{rejected} as real sharing with the holder named, and the message offers
\code{Dict.update}, which is accepted (Appendix~\ref{app:groupby}). Lean accepts the lookup form and
copies silently at run time, because the dictionary still holds the group; the program is
quadratic \citep[\S2.4]{huisinga2023lean}. \emph{Solved with the same cost.}

\subsection{Recursion and folds}
\label{sec:hard-recursion}

\begin{Verbatim}
go xs acc = case xs of
    [] → acc
    [ x, …rest ] → go rest (List.push acc (f x))
\end{Verbatim}
\code{go}'s group is one declaration. The optimistic start assumes \code{acc} borrowed and the result
fresh. In round 1 the tail call passes \code{List.push acc …}, a demand on \code{acc}'s cell, whose only
holder is \code{acc} itself; after the jump \code{acc} is rebound, so it is dead. Hence
$\mathit{use}(\mathit{acc}) = \consumed$, and $\mathit{res} = \{\pi_2\}$ (the \code{[]} arm) united with
the recursive call's result, $\{\fresh\}$ under the round's assumption. Round 2, with
$\mathit{res} = \{\pi_2, \fresh\}$, gives the same; validation passes. \code{List.foldl} has this shape,
with its callback's class carrying the accumulator's ownership, so
\code{List.foldl xs [] (λx acc → List.push acc x)} is accepted and runs as the in-place builds of our
prototype study did. Mutual recursion is the same fixpoint over the group. \emph{Solved.}

\subsection{Branches that keep the old value}
\label{sec:hard-branches}

\begin{Verbatim}
case List.get xs i of
    Just v  → v
    Nothing → List.set xs i d ▷ List.get …
\end{Verbatim}
\code{v} is an element, not the array, and per-arm liveness makes \code{xs} dead after the \code{set} on
the \code{Nothing} arm. Accepted. The shape that Rust's borrow checker long rejected (the third
``problem case'' of non-lexical lifetimes \citep{rfc2094}) does not arise, because beni has no
references into a list: \code{get} returns the element, not a pointer to the slot.

In \code{ys = if c then List.set xs 0 1 else xs}, $A(\mathit{ys})$ is \code{xs}'s cell on both arms;
the demand in the \code{then} arm asks whether \code{xs} is live after it \emph{on that arm}, and it is
not. Accepted, and \code{ys} owns the cell. Futhark's rule that the aliases of a conditional are the
union of both branches \citep[Fig.~5]{henriksen2017futhark} is kept for cells, but liveness is per path,
so a value ``used'' only in the arm not taken causes no false error. \emph{Solved.}

What remains coarse: in \code{xs = if c then a else b} followed by \code{List.set xs 0 1}, both \code{a}
and \code{b} must be unique, so a later read of \code{a} on the path where \code{xs} is \code{b} is a
false error. The evaluation counts how often this occurs.

\subsection{The foreign boundary}
\label{sec:hard-foreign}

A \code{foreign} has no body; its declaration carries a summary as it already carries an effect rung.
\emph{The default is pessimistic}: a parameter that says nothing is shared (beni's existing contract
lets a sibling keep a list it was given in order to read it later), and a result that says nothing is
$\top$. A sibling that only reads its argument during the call says borrowed; one that returns a fresh
array that nothing else holds says fresh; \code{Debug.log}'s result says that it is its second
argument. A platform author who writes nothing therefore loses precision and never soundness.
\code{Js.from x} marks \code{x} escaped forever; a \code{Js.call} that passes a list is a call with no
summary and escapes it. The residual risk is a sibling that lies---a bug in privileged code of the same
class as a wrong effect rung. \emph{Solved, by contract} (A2).

\subsection{Fibers and commands that capture values}
\label{sec:hard-fibers}

\code{Cmd.perform (λsend → send (Saved model.todos))} captures \code{model.todos} in a thunk that the
runtime stores and runs later. \code{Cmd.perform}'s argument has $\mathit{escapes} = \mathit{yes}$, so
every capture of the lambda is escaped at the capture, and \code{model.todos}'s cell is shared from
then on, in this arm \emph{and in every later dispatch that receives the same cell in its model}
(\S\ref{sec:hard-tea}). Spawning a fiber, a parallel thunk, \code{Ref.set} and \code{Queue.offer} are
the same. A thunk the callee calls synchronously and does not keep---the body of a structured
concurrency scope, the acquire step of a bracket---has $\mathit{calls} \le 1$,
$\mathit{escapes} = \mathit{no}$, and captures nothing permanently; a bracket's \emph{release} is stored
on the fiber's finaliser list and escapes. \emph{Solved by escape}, with the cost that a list handed to
any command or fiber is shared until the end of the program, because the analysis cannot see when the
fiber reads it.

\subsection{The TEA update: a hand-over protocol}
\label{sec:hard-tea}

beni's current template runtime cannot host the rule. It keeps the rows it rendered and the last
value of every hole, so every rendered list is a content holder the program cannot see, and a
\code{List.push model.rows x} in \code{update} would be an error on every message. Our prototype study
measured the same thing from the other side: with the runtime as written, no scheme wrote the table
in place. The direct platform (\S\ref{sec:bg-beni}) removes these holds by design: there is no
\code{view} at run time and no retained render, slots hold leaf values, instance arrays hold row
items and compare them by identity, and conditional DOM edits are guarded by their own tags, never by
list identity. What remains is to state the contract the checker reads, as six obligations on the
platform's lowering and runtime.

\begin{description}[leftmargin=1.2em,font=\normalfont\itshape]
\item[Hand-over.] The handler calls the arm with the model \emph{consumed}: after the arm, the runtime
holds no content reference to any list reachable from the old model. The module-level model variable
is reassigned by the arm.
\item[Reads before the arm.] Anything the handler needs from the old model is read \emph{before} the
arm, into locals. A list read this way is an ordinary holder of the handler, analysed by the same
rules; a handler that read the old list and compared it after the arm would commit a sharing error,
which the lowering must not generate.
\item[Item reads.] Row items in instance arrays are read only by the lowering's own code: compared by
\code{===} (A4), with a row's DOM writes reading the \emph{new} item. The one program-visible read
through an instance is a \emph{listener body}, which reads its handler's arguments when the event
fires---\code{onClick=\{Select model.id\}} reads \code{model.id} at the click. A listener body is
therefore program code the checker analyses; it reads the model and the item \emph{before} the arm.
\item[Slots.] A derived value's slot that holds a list is recomputed by every message whose write set
conflicts with its reads before any DOM write reads it, and the old slot value is never dereferenced
after the arm. The write-set analysis's soundness already covers this once ``changed'' is read as
``possibly changed in content'' rather than ``not \code{===}''.
\item[Entry fixpoint.] \code{update} has no caller to discharge \Unescaped, \Dead\ and \Disjoint\ for
its model parameter, so the checker computes the model's abstract value at entry itself: the join, over
\code{init}'s result and every arm's result, of the cell sets at every model path, iterated over the
arms to a fixpoint (one round per cell the arms can place; in practice two). Two facts come out of it.
\emph{Aliasing between paths}: \code{init} of the form \code{xs = List.range 1 3; \{ a = xs, b = xs \}}
places one cell at $.a$ and $.b$, so a push on \code{model.a} fails \Disjoint\ against
$(\mathit{model}, .b)$. \emph{Escapes across dispatches}: a cell that any arm captures in a command,
a subscription, a stored closure, a \code{Debug} or \code{Js} call, a send, or an event payload (below)
is escaped at every entry, because the runtime may hand the same cell back in the next model. The
entry summary is then $\mathit{use}(\mathit{model}, p) = \shared$ for every escaped or aliased path,
and $\consumed$ otherwise. This is the one whole-program computation in the design.
\item[Payloads.] A message's payload list is unique at the arm \emph{only if nothing else holds it when
the arm runs}. A program sender---a fiber calling \code{send (Got rows)}---gives it up: sending consumes
the message, so a sender that reads \code{rows} afterwards is rejected at the sender. A \emph{view
event} whose handler expression is a model path or a row item (\code{onClick=\{Keep model.rows\}})
hands the arm a cell the model still holds; such a payload path is shared at the arm, and the message
offers to copy it in the arm or to send an identifier instead of the list.
\end{description}

Under these obligations, the derivation for a table application's row swap and row update, and for
\code{List.update model.rows 1 …} and \code{List.push model.rows …} in general, is: the model is consumed
at entry; \code{model.rows} is not escaped by any arm; the operand's only holder is
$(\mathit{model}, .\mathit{rows})$; the model is dead after the arm. The update is accepted, and the
write is \code{a[i] = v} on a plain array, derived from the general rule rather than from a
model-specific one. \emph{Solved, conditional on the platform keeping the six obligations}; the entry
fixpoint is real work.

\paragraph{Derived values are reads, never demands.} The direct platform compiles a \code{let} in
\code{view} into a derived value, recomputed by handlers while the page holds the model; a TodoMVC
\code{<For each=\{List.filter model.todos (visible model.filter \_)\}>} is one. A development-mode check
re-evaluates every derived value after each dispatch and expects no change. No expression that a
derived value or that check evaluates may therefore be a demand. Under the nine-writer demand set this
is automatic, since \code{filter} and \code{map} are not demands; if the demand set is widened it must be
stated, and it is the reason the widening cannot be a single function (\S\ref{sec:disc-options}).

\subsection{The explicit copy}
\label{sec:hard-copy}

\code{List.copy : List a [borrowed] → List a [fresh]} is a shallow copy, \code{a.slice()}. The elements
are the same values; records are immutable and need no copying; a list nested inside an element is
still shared, and editing it is still an error whose message names the inner holder. A copy has a new
identity, so a keyed list in the view sees new rows, which is right for a list whose contents the
program is about to change behind the original's back. Copying a list that is already unique is legal
and wasteful; the evaluation counts such copies to decide whether a warning is wanted---a warning, never
an error. \emph{Solved.}

\subsection{Views}
\label{sec:hard-views}

\code{[ x, …rest ]} binds \code{rest} to a view of \code{xs}'s base array. A view is a holder of the base,
and a write through a view writes the base: \code{List.set rest i v} is \code{b[o + i] = v}. The demand
on \code{rest} therefore asks that every holder of the base---\code{xs} included---be dead. In the common
walk (\code{go rest …} with \code{xs} not read again) it is; where \code{xs} is read afterwards, the
rejection is correct, because the elements \emph{are} shared. A view is always a suffix of its base, so
\code{push} on a view is a push on the base and a new header, in $O(1)$. \emph{Solved.}

\subsection{Mutable cells}
\label{sec:hard-ref}

\code{Ref.get r} returns an alias of whatever the cell holds, and \code{Ref.set r xs} stores \code{xs};
the analysis cannot see who else has called \code{Ref.get}. A list stored in a \code{Ref} is therefore
escaped, every \code{Ref.get} result is $\top$-held, and \code{Ref.update r (λxs → List.push xs x)} is
rejected as a limit. A linear-cell discipline would need a proof that no \code{Ref.get} result is live
across the update, a whole-program fact about an unbounded set of readers. The fix is a copy in the
callback, or a persistent structure in the cell. \emph{Open, as a precision limit.}

\subsection{Messages as values}
\label{sec:hard-messages}

A message stored in the model, logged whole, or mapped by \code{Html.map} with a function that is not a
constructor is a value; its payload lists are shared wherever the message is. A payload passed as
handler parameters---the common case---is covered by the payload obligation. \emph{Solved.}

\subsection{Versions: undo histories, diffs, snapshots}
\label{sec:hard-undo}

In \code{\{ model | history = [ model.todos, …model.history ], todos = List.push model.todos t \}} the
record update evaluates its fields in written order, so \code{history} holds the old array when the
push runs---a live holder. The program is \emph{rejected}, correctly: the undo entry would change with
the push. The message offers \code{List.copy model.todos} in the history, $O(n)$ per edit, the price of a
version. A persistent vector written in beni as a library module would give $O(\log n)$ versions to
programs that want them; it would be a library, not the representation of \code{List}. \emph{Solved},
with the cost stated.

\subsection{Debugging aids}
\label{sec:hard-debug}

\code{Debug.log tag x} reads \code{x} at the call and returns it; \code{Debug.toString} reads at the call.
Neither is a holder afterwards. A development crash screen must print only the exception and the
stack (A9); one that printed the model would read a half-updated record, a difference in a debugging
aid after the program has stopped, not in the program. \emph{Solved}, with A9 recorded.

\subsection{Inlining and specialisation}
\label{sec:hard-optimiser}

Inlining substitutes a body at a call, which preserves evaluation order and holders. The verdict was
computed on the source and is not recomputed. Whole-program specialisation may copy a callee per call
site, and a copy may be more unique than the general summary; the rule ignores this, so the accepted
set is the source's, never the optimiser's. \emph{Solved by construction}, at the cost that a program
the specialiser could have proved safe is still rejected.
